
\documentclass[10pt]{article} % For LaTeX2e
%\usepackage{tmlr}
% If accepted, instead use the following line for the camera-ready submission:
%\usepackage[accepted]{tmlr}
% To de-anonymize and remove mentions to TMLR (for example for posting to preprint servers), instead use the following:
\usepackage[preprint]{tmlr}

% Optional math commands from https://github.com/goodfeli/dlbook_notation.
%%%%% NEW MATH DEFINITIONS %%%%%

\usepackage{amsmath,amsfonts,bm}

% Mark sections of captions for referring to divisions of figures
\newcommand{\figleft}{{\em (Left)}}
\newcommand{\figcenter}{{\em (Center)}}
\newcommand{\figright}{{\em (Right)}}
\newcommand{\figtop}{{\em (Top)}}
\newcommand{\figbottom}{{\em (Bottom)}}
\newcommand{\captiona}{{\em (a)}}
\newcommand{\captionb}{{\em (b)}}
\newcommand{\captionc}{{\em (c)}}
\newcommand{\captiond}{{\em (d)}}

% Highlight a newly defined term
\newcommand{\newterm}[1]{{\bf #1}}


% Figure reference, lower-case.
\def\figref#1{figure~\ref{#1}}
% Figure reference, capital. For start of sentence
\def\Figref#1{Figure~\ref{#1}}
\def\twofigref#1#2{figures \ref{#1} and \ref{#2}}
\def\quadfigref#1#2#3#4{figures \ref{#1}, \ref{#2}, \ref{#3} and \ref{#4}}
% Section reference, lower-case.
\def\secref#1{section~\ref{#1}}
% Section reference, capital.
\def\Secref#1{Section~\ref{#1}}
% Reference to two sections.
\def\twosecrefs#1#2{sections \ref{#1} and \ref{#2}}
% Reference to three sections.
\def\secrefs#1#2#3{sections \ref{#1}, \ref{#2} and \ref{#3}}
% Reference to an equation, lower-case.
\def\eqref#1{equation~\ref{#1}}
% Reference to an equation, upper case
\def\Eqref#1{Equation~\ref{#1}}
% A raw reference to an equation---avoid using if possible
\def\plaineqref#1{\ref{#1}}
% Reference to a chapter, lower-case.
\def\chapref#1{chapter~\ref{#1}}
% Reference to an equation, upper case.
\def\Chapref#1{Chapter~\ref{#1}}
% Reference to a range of chapters
\def\rangechapref#1#2{chapters\ref{#1}--\ref{#2}}
% Reference to an algorithm, lower-case.
\def\algref#1{algorithm~\ref{#1}}
% Reference to an algorithm, upper case.
\def\Algref#1{Algorithm~\ref{#1}}
\def\twoalgref#1#2{algorithms \ref{#1} and \ref{#2}}
\def\Twoalgref#1#2{Algorithms \ref{#1} and \ref{#2}}
% Reference to a part, lower case
\def\partref#1{part~\ref{#1}}
% Reference to a part, upper case
\def\Partref#1{Part~\ref{#1}}
\def\twopartref#1#2{parts \ref{#1} and \ref{#2}}

\def\ceil#1{\lceil #1 \rceil}
\def\floor#1{\lfloor #1 \rfloor}
\def\1{\bm{1}}
\newcommand{\train}{\mathcal{D}}
\newcommand{\valid}{\mathcal{D_{\mathrm{valid}}}}
\newcommand{\test}{\mathcal{D_{\mathrm{test}}}}

\def\eps{{\epsilon}}


% Random variables
\def\reta{{\textnormal{$\eta$}}}
\def\ra{{\textnormal{a}}}
\def\rb{{\textnormal{b}}}
\def\rc{{\textnormal{c}}}
\def\rd{{\textnormal{d}}}
\def\re{{\textnormal{e}}}
\def\rf{{\textnormal{f}}}
\def\rg{{\textnormal{g}}}
\def\rh{{\textnormal{h}}}
\def\ri{{\textnormal{i}}}
\def\rj{{\textnormal{j}}}
\def\rk{{\textnormal{k}}}
\def\rl{{\textnormal{l}}}
% rm is already a command, just don't name any random variables m
\def\rn{{\textnormal{n}}}
\def\ro{{\textnormal{o}}}
\def\rp{{\textnormal{p}}}
\def\rq{{\textnormal{q}}}
\def\rr{{\textnormal{r}}}
\def\rs{{\textnormal{s}}}
\def\rt{{\textnormal{t}}}
\def\ru{{\textnormal{u}}}
\def\rv{{\textnormal{v}}}
\def\rw{{\textnormal{w}}}
\def\rx{{\textnormal{x}}}
\def\ry{{\textnormal{y}}}
\def\rz{{\textnormal{z}}}

% Random vectors
\def\rvepsilon{{\mathbf{\epsilon}}}
\def\rvtheta{{\mathbf{\theta}}}
\def\rva{{\mathbf{a}}}
\def\rvb{{\mathbf{b}}}
\def\rvc{{\mathbf{c}}}
\def\rvd{{\mathbf{d}}}
\def\rve{{\mathbf{e}}}
\def\rvf{{\mathbf{f}}}
\def\rvg{{\mathbf{g}}}
\def\rvh{{\mathbf{h}}}
\def\rvu{{\mathbf{i}}}
\def\rvj{{\mathbf{j}}}
\def\rvk{{\mathbf{k}}}
\def\rvl{{\mathbf{l}}}
\def\rvm{{\mathbf{m}}}
\def\rvn{{\mathbf{n}}}
\def\rvo{{\mathbf{o}}}
\def\rvp{{\mathbf{p}}}
\def\rvq{{\mathbf{q}}}
\def\rvr{{\mathbf{r}}}
\def\rvs{{\mathbf{s}}}
\def\rvt{{\mathbf{t}}}
\def\rvu{{\mathbf{u}}}
\def\rvv{{\mathbf{v}}}
\def\rvw{{\mathbf{w}}}
\def\rvx{{\mathbf{x}}}
\def\rvy{{\mathbf{y}}}
\def\rvz{{\mathbf{z}}}

% Elements of random vectors
\def\erva{{\textnormal{a}}}
\def\ervb{{\textnormal{b}}}
\def\ervc{{\textnormal{c}}}
\def\ervd{{\textnormal{d}}}
\def\erve{{\textnormal{e}}}
\def\ervf{{\textnormal{f}}}
\def\ervg{{\textnormal{g}}}
\def\ervh{{\textnormal{h}}}
\def\ervi{{\textnormal{i}}}
\def\ervj{{\textnormal{j}}}
\def\ervk{{\textnormal{k}}}
\def\ervl{{\textnormal{l}}}
\def\ervm{{\textnormal{m}}}
\def\ervn{{\textnormal{n}}}
\def\ervo{{\textnormal{o}}}
\def\ervp{{\textnormal{p}}}
\def\ervq{{\textnormal{q}}}
\def\ervr{{\textnormal{r}}}
\def\ervs{{\textnormal{s}}}
\def\ervt{{\textnormal{t}}}
\def\ervu{{\textnormal{u}}}
\def\ervv{{\textnormal{v}}}
\def\ervw{{\textnormal{w}}}
\def\ervx{{\textnormal{x}}}
\def\ervy{{\textnormal{y}}}
\def\ervz{{\textnormal{z}}}

% Random matrices
\def\rmA{{\mathbf{A}}}
\def\rmB{{\mathbf{B}}}
\def\rmC{{\mathbf{C}}}
\def\rmD{{\mathbf{D}}}
\def\rmE{{\mathbf{E}}}
\def\rmF{{\mathbf{F}}}
\def\rmG{{\mathbf{G}}}
\def\rmH{{\mathbf{H}}}
\def\rmI{{\mathbf{I}}}
\def\rmJ{{\mathbf{J}}}
\def\rmK{{\mathbf{K}}}
\def\rmL{{\mathbf{L}}}
\def\rmM{{\mathbf{M}}}
\def\rmN{{\mathbf{N}}}
\def\rmO{{\mathbf{O}}}
\def\rmP{{\mathbf{P}}}
\def\rmQ{{\mathbf{Q}}}
\def\rmR{{\mathbf{R}}}
\def\rmS{{\mathbf{S}}}
\def\rmT{{\mathbf{T}}}
\def\rmU{{\mathbf{U}}}
\def\rmV{{\mathbf{V}}}
\def\rmW{{\mathbf{W}}}
\def\rmX{{\mathbf{X}}}
\def\rmY{{\mathbf{Y}}}
\def\rmZ{{\mathbf{Z}}}

% Elements of random matrices
\def\ermA{{\textnormal{A}}}
\def\ermB{{\textnormal{B}}}
\def\ermC{{\textnormal{C}}}
\def\ermD{{\textnormal{D}}}
\def\ermE{{\textnormal{E}}}
\def\ermF{{\textnormal{F}}}
\def\ermG{{\textnormal{G}}}
\def\ermH{{\textnormal{H}}}
\def\ermI{{\textnormal{I}}}
\def\ermJ{{\textnormal{J}}}
\def\ermK{{\textnormal{K}}}
\def\ermL{{\textnormal{L}}}
\def\ermM{{\textnormal{M}}}
\def\ermN{{\textnormal{N}}}
\def\ermO{{\textnormal{O}}}
\def\ermP{{\textnormal{P}}}
\def\ermQ{{\textnormal{Q}}}
\def\ermR{{\textnormal{R}}}
\def\ermS{{\textnormal{S}}}
\def\ermT{{\textnormal{T}}}
\def\ermU{{\textnormal{U}}}
\def\ermV{{\textnormal{V}}}
\def\ermW{{\textnormal{W}}}
\def\ermX{{\textnormal{X}}}
\def\ermY{{\textnormal{Y}}}
\def\ermZ{{\textnormal{Z}}}

% Vectors
\def\vzero{{\bm{0}}}
\def\vone{{\bm{1}}}
\def\vmu{{\bm{\mu}}}
\def\vtheta{{\bm{\theta}}}
\def\va{{\bm{a}}}
\def\vb{{\bm{b}}}
\def\vc{{\bm{c}}}
\def\vd{{\bm{d}}}
\def\ve{{\bm{e}}}
\def\vf{{\bm{f}}}
\def\vg{{\bm{g}}}
\def\vh{{\bm{h}}}
\def\vi{{\bm{i}}}
\def\vj{{\bm{j}}}
\def\vk{{\bm{k}}}
\def\vl{{\bm{l}}}
\def\vm{{\bm{m}}}
\def\vn{{\bm{n}}}
\def\vo{{\bm{o}}}
\def\vp{{\bm{p}}}
\def\vq{{\bm{q}}}
\def\vr{{\bm{r}}}
\def\vs{{\bm{s}}}
\def\vt{{\bm{t}}}
\def\vu{{\bm{u}}}
\def\vv{{\bm{v}}}
\def\vw{{\bm{w}}}
\def\vx{{\bm{x}}}
\def\vy{{\bm{y}}}
\def\vz{{\bm{z}}}

% Elements of vectors
\def\evalpha{{\alpha}}
\def\evbeta{{\beta}}
\def\evepsilon{{\epsilon}}
\def\evlambda{{\lambda}}
\def\evomega{{\omega}}
\def\evmu{{\mu}}
\def\evpsi{{\psi}}
\def\evsigma{{\sigma}}
\def\evtheta{{\theta}}
\def\eva{{a}}
\def\evb{{b}}
\def\evc{{c}}
\def\evd{{d}}
\def\eve{{e}}
\def\evf{{f}}
\def\evg{{g}}
\def\evh{{h}}
\def\evi{{i}}
\def\evj{{j}}
\def\evk{{k}}
\def\evl{{l}}
\def\evm{{m}}
\def\evn{{n}}
\def\evo{{o}}
\def\evp{{p}}
\def\evq{{q}}
\def\evr{{r}}
\def\evs{{s}}
\def\evt{{t}}
\def\evu{{u}}
\def\evv{{v}}
\def\evw{{w}}
\def\evx{{x}}
\def\evy{{y}}
\def\evz{{z}}

% Matrix
\def\mA{{\bm{A}}}
\def\mB{{\bm{B}}}
\def\mC{{\bm{C}}}
\def\mD{{\bm{D}}}
\def\mE{{\bm{E}}}
\def\mF{{\bm{F}}}
\def\mG{{\bm{G}}}
\def\mH{{\bm{H}}}
\def\mI{{\bm{I}}}
\def\mJ{{\bm{J}}}
\def\mK{{\bm{K}}}
\def\mL{{\bm{L}}}
\def\mM{{\bm{M}}}
\def\mN{{\bm{N}}}
\def\mO{{\bm{O}}}
\def\mP{{\bm{P}}}
\def\mQ{{\bm{Q}}}
\def\mR{{\bm{R}}}
\def\mS{{\bm{S}}}
\def\mT{{\bm{T}}}
\def\mU{{\bm{U}}}
\def\mV{{\bm{V}}}
\def\mW{{\bm{W}}}
\def\mX{{\bm{X}}}
\def\mY{{\bm{Y}}}
\def\mZ{{\bm{Z}}}
\def\mBeta{{\bm{\beta}}}
\def\mPhi{{\bm{\Phi}}}
\def\mLambda{{\bm{\Lambda}}}
\def\mSigma{{\bm{\Sigma}}}

% Tensor
\DeclareMathAlphabet{\mathsfit}{\encodingdefault}{\sfdefault}{m}{sl}
\SetMathAlphabet{\mathsfit}{bold}{\encodingdefault}{\sfdefault}{bx}{n}
\newcommand{\tens}[1]{\bm{\mathsfit{#1}}}
\def\tA{{\tens{A}}}
\def\tB{{\tens{B}}}
\def\tC{{\tens{C}}}
\def\tD{{\tens{D}}}
\def\tE{{\tens{E}}}
\def\tF{{\tens{F}}}
\def\tG{{\tens{G}}}
\def\tH{{\tens{H}}}
\def\tI{{\tens{I}}}
\def\tJ{{\tens{J}}}
\def\tK{{\tens{K}}}
\def\tL{{\tens{L}}}
\def\tM{{\tens{M}}}
\def\tN{{\tens{N}}}
\def\tO{{\tens{O}}}
\def\tP{{\tens{P}}}
\def\tQ{{\tens{Q}}}
\def\tR{{\tens{R}}}
\def\tS{{\tens{S}}}
\def\tT{{\tens{T}}}
\def\tU{{\tens{U}}}
\def\tV{{\tens{V}}}
\def\tW{{\tens{W}}}
\def\tX{{\tens{X}}}
\def\tY{{\tens{Y}}}
\def\tZ{{\tens{Z}}}


% Graph
\def\gA{{\mathcal{A}}}
\def\gB{{\mathcal{B}}}
\def\gC{{\mathcal{C}}}
\def\gD{{\mathcal{D}}}
\def\gE{{\mathcal{E}}}
\def\gF{{\mathcal{F}}}
\def\gG{{\mathcal{G}}}
\def\gH{{\mathcal{H}}}
\def\gI{{\mathcal{I}}}
\def\gJ{{\mathcal{J}}}
\def\gK{{\mathcal{K}}}
\def\gL{{\mathcal{L}}}
\def\gM{{\mathcal{M}}}
\def\gN{{\mathcal{N}}}
\def\gO{{\mathcal{O}}}
\def\gP{{\mathcal{P}}}
\def\gQ{{\mathcal{Q}}}
\def\gR{{\mathcal{R}}}
\def\gS{{\mathcal{S}}}
\def\gT{{\mathcal{T}}}
\def\gU{{\mathcal{U}}}
\def\gV{{\mathcal{V}}}
\def\gW{{\mathcal{W}}}
\def\gX{{\mathcal{X}}}
\def\gY{{\mathcal{Y}}}
\def\gZ{{\mathcal{Z}}}

% Sets
\def\sA{{\mathbb{A}}}
\def\sB{{\mathbb{B}}}
\def\sC{{\mathbb{C}}}
\def\sD{{\mathbb{D}}}
% Don't use a set called E, because this would be the same as our symbol
% for expectation.
\def\sF{{\mathbb{F}}}
\def\sG{{\mathbb{G}}}
\def\sH{{\mathbb{H}}}
\def\sI{{\mathbb{I}}}
\def\sJ{{\mathbb{J}}}
\def\sK{{\mathbb{K}}}
\def\sL{{\mathbb{L}}}
\def\sM{{\mathbb{M}}}
\def\sN{{\mathbb{N}}}
\def\sO{{\mathbb{O}}}
\def\sP{{\mathbb{P}}}
\def\sQ{{\mathbb{Q}}}
\def\sR{{\mathbb{R}}}
\def\sS{{\mathbb{S}}}
\def\sT{{\mathbb{T}}}
\def\sU{{\mathbb{U}}}
\def\sV{{\mathbb{V}}}
\def\sW{{\mathbb{W}}}
\def\sX{{\mathbb{X}}}
\def\sY{{\mathbb{Y}}}
\def\sZ{{\mathbb{Z}}}

% Entries of a matrix
\def\emLambda{{\Lambda}}
\def\emA{{A}}
\def\emB{{B}}
\def\emC{{C}}
\def\emD{{D}}
\def\emE{{E}}
\def\emF{{F}}
\def\emG{{G}}
\def\emH{{H}}
\def\emI{{I}}
\def\emJ{{J}}
\def\emK{{K}}
\def\emL{{L}}
\def\emM{{M}}
\def\emN{{N}}
\def\emO{{O}}
\def\emP{{P}}
\def\emQ{{Q}}
\def\emR{{R}}
\def\emS{{S}}
\def\emT{{T}}
\def\emU{{U}}
\def\emV{{V}}
\def\emW{{W}}
\def\emX{{X}}
\def\emY{{Y}}
\def\emZ{{Z}}
\def\emSigma{{\Sigma}}

% entries of a tensor
% Same font as tensor, without \bm wrapper
\newcommand{\etens}[1]{\mathsfit{#1}}
\def\etLambda{{\etens{\Lambda}}}
\def\etA{{\etens{A}}}
\def\etB{{\etens{B}}}
\def\etC{{\etens{C}}}
\def\etD{{\etens{D}}}
\def\etE{{\etens{E}}}
\def\etF{{\etens{F}}}
\def\etG{{\etens{G}}}
\def\etH{{\etens{H}}}
\def\etI{{\etens{I}}}
\def\etJ{{\etens{J}}}
\def\etK{{\etens{K}}}
\def\etL{{\etens{L}}}
\def\etM{{\etens{M}}}
\def\etN{{\etens{N}}}
\def\etO{{\etens{O}}}
\def\etP{{\etens{P}}}
\def\etQ{{\etens{Q}}}
\def\etR{{\etens{R}}}
\def\etS{{\etens{S}}}
\def\etT{{\etens{T}}}
\def\etU{{\etens{U}}}
\def\etV{{\etens{V}}}
\def\etW{{\etens{W}}}
\def\etX{{\etens{X}}}
\def\etY{{\etens{Y}}}
\def\etZ{{\etens{Z}}}

% The true underlying data generating distribution
\newcommand{\pdata}{p_{\rm{data}}}
% The empirical distribution defined by the training set
\newcommand{\ptrain}{\hat{p}_{\rm{data}}}
\newcommand{\Ptrain}{\hat{P}_{\rm{data}}}
% The model distribution
\newcommand{\pmodel}{p_{\rm{model}}}
\newcommand{\Pmodel}{P_{\rm{model}}}
\newcommand{\ptildemodel}{\tilde{p}_{\rm{model}}}
% Stochastic autoencoder distributions
\newcommand{\pencode}{p_{\rm{encoder}}}
\newcommand{\pdecode}{p_{\rm{decoder}}}
\newcommand{\precons}{p_{\rm{reconstruct}}}

\newcommand{\laplace}{\mathrm{Laplace}} % Laplace distribution

\newcommand{\E}{\mathbb{E}}
\newcommand{\Ls}{\mathcal{L}}
\newcommand{\R}{\mathbb{R}}
\newcommand{\emp}{\tilde{p}}
\newcommand{\lr}{\alpha}
\newcommand{\reg}{\lambda}
\newcommand{\rect}{\mathrm{rectifier}}
\newcommand{\softmax}{\mathrm{softmax}}
\newcommand{\sigmoid}{\sigma}
\newcommand{\softplus}{\zeta}
\newcommand{\KL}{D_{\mathrm{KL}}}
\newcommand{\Var}{\mathrm{Var}}
\newcommand{\standarderror}{\mathrm{SE}}
\newcommand{\Cov}{\mathrm{Cov}}
% Wolfram Mathworld says $L^2$ is for function spaces and $\ell^2$ is for vectors
% But then they seem to use $L^2$ for vectors throughout the site, and so does
% wikipedia.
\newcommand{\normlzero}{L^0}
\newcommand{\normlone}{L^1}
\newcommand{\normltwo}{L^2}
\newcommand{\normlp}{L^p}
\newcommand{\normmax}{L^\infty}

\newcommand{\parents}{Pa} % See usage in notation.tex. Chosen to match Daphne's book.

\DeclareMathOperator*{\argmax}{arg\,max}
\DeclareMathOperator*{\argmin}{arg\,min}

\DeclareMathOperator{\sign}{sign}
\DeclareMathOperator{\Tr}{Tr}
\let\ab\allowbreak


\usepackage{hyperref}
\usepackage{url}


\title{On the Reliability of Neural Networks}
% [Eng] A formal methodology for the Reliability of Neural Networks
% Authors must not appear in the submitted version. They should be hidden
% as long as the tmlr package is used without the [accepted] or [preprint] options.
% Non-anonymous submissions will be rejected without review.

\author{\name Rosario Pio Gnazzo \email rosariopiognazzo@gmail.com \\
      \addr Department of Computer Science\\
      University of Salerno}

% The \author macro works with any number of authors. Use \AND 
% to separate the names and addresses of multiple authors.

\newcommand{\fix}{\marginpar{FIX}}
\newcommand{\new}{\marginpar{NEW}}

\def\month{MM}  % Insert correct month for camera-ready version
\def\year{YYYY} % Insert correct year for camera-ready version
\def\openreview{\url{https://openreview.net/forum?id=XXXX}} % Insert correct link to OpenReview for camera-ready version


\begin{document}


\maketitle

\begin{abstract}
Although Deep Learning has achieved excellent performance in numerous domains, it lacks \textit{formal metrics} to quantify the safety margin and the predictive robustness in \textit{safety-critical} contexts. Normally the focus is on evaluating a model solely on evaluation metrics (performance-based), without taking into account a view on the intrinsic \textit{reliability} of the system.

This work attempts to study a mathematical-statistical method of analysis by integrating \textit{classical Reliability Theory} into \textit{Deep Learning}. The idea is, by modeling neural layers as stochastic systems of the \textit{k-out-of-n} type, to focus on \textit{neuronal activations} rather than on weights, using the Beta-Binomial distribution to represent the behavior (understood as the probability of being \textit{active} when the network achieves a success) of the neurons.

The method makes it possible to estimate \textit{point reliability parameters}, minimum \textit{activation thresholds} $s_{(k)}$ and \textit{operational metrics} $\rho_{(k)}$. Moreover, in order to handle the uncertainty of the estimates and to have a more robust analysis, the \textit{confidence intervals} are derived analytically.
\end{abstract}

\section{Introduction}
\label{intro}
One of the problems encountered in deep learning is that it is not possible to obtain a robust prediction of how a model will behave on test data. Normally the focus is on evaluating the \textit{accuracy} of the model (or other metrics), without being aware of a \textit{safety margin} of the network, that is: \textit{how reliable it is in carrying out the task}. In this study I will present a model that proposes exactly this, exploiting the ideas already known in the field of the reliability of multi-component systems\footnote{Reliability theory has been studied for years in engineering and computer science as a tool to evaluate a system on the basis of the properties of the components it is made of. The intuition from which the work starts is that of considering a neural network as a system of components that, connected to each other, have statistical properties describing the reliability of the network.}. 
The approach focuses on neuronal activations, treating individual neurons as if they were physical components of a system.

\section{Definitions}
\label{assunzioni}
Before focusing on the approach, we define the objects that compose it and their characteristics.

\subsection{Neural Network}
\label{A1-FeedForward}
We denote with the notation $f: X^{n} \xrightarrow{\mathbb{H}^{L}} Y^{m}$ a neural network such that:
\begin{itemize}
    \item $X^{n} = \{x_{1}, x_{2}, \dots, x_{n}\}$, identifies the \textit{input layer} composed of $n$ neurons;
    \item $\mathbb{H}^{L} = \{H^{(1)}_{n_{1}}, H^{(2)}_{n_{2}}, \dots, H^{(L)}_{n_{L}}\}$ which identifies the \textit{set of the} $L$ \textit{internal layers} each composed of $n_{l}$ neurons;
    \begin{equation}
        H^{(l)}_{n_{l}} = \{h^{(l)}_{1}, h^{(l)}_{2}, \dots, h^{(l)}_{n_{l}}\}\qquad l = 1, \dots, L
    \end{equation}
    \item $Y^{m} = \{y_{1}, y_{2}, \dots, y_{m}\}$, identifies the \textit{output layer} composed of $m$ neurons.
\end{itemize}
Given this formalization we define
\begin{itemize}
    \item $z^{(k)}_{i}$ the \textit{Activation of the} \textit{i-th neuron} belonging to the \textit{k-th} layer after the application of the activation function $\sigma$;
\end{itemize}

\subsection{Neuronal activation}
\label{A2-Attivazione}
In order to understand the definition of reliability that will be given later in the work, it is necessary to define what the \textit{activation} state of a neuron represents in probabilistic terms. With $A^{(k)}_{i}$ we indicate that neuron \textit{i} of the \textit{k-th} layer is active. This definition can be formalized as a binary function
\begin{equation}
    A^{(k)}_{i} = \mathbb{I}[z^{(k)}_{i} > \tau^{(l)}] = 
    \begin{cases} 
        1, & \mbox{if }z^{(k)}_{i} > \tau^{(l)} \\ 
        0, & \mbox{otherwise}
    \end{cases}
\end{equation}
where $\tau^{(l)}$ denotes the activation threshold, which depends on the layer.\footnote{e.g. If $\tau^{(l)} = 0$ then the activation function we are considering is \textit{ReLU}}.

\subsection{System component}
\label{A3-Component}
From the definition [\ref{A2-Attivazione}] we can describe a neuron as a component of a system observed under its own activation condition. Given a neuron $n^{(k)}_{i}$ belonging to any layer \textit{k} of a neural network [\ref{A1-FeedForward}], we define the event $B = \{\text{Observing a \textit{correct} output from the network}\}$, then the component $C^{(k)}_{i}$ related to neuron $n^{(k)}_{i}$ will be described by a \textit{Bernoulli r.v.}
\begin{equation} 
    C^{(k)}_{i} \sim Bernoulli(R^{(k)}_{i})
\end{equation}
where $R^{(k)}_{i}$\footnote{Empirically it is computed using the frequentist definition of probability $\frac{\mbox{No. of Activations with Correct Output}}{\mbox{Tot. No. of Correct Outputs}}$}
measures the \textit{probability of participation in the success}, that is, the probability that neuron \textit{i} turns on when the network produces a correct result.
\begin{equation}
    R^{(k)}_{i} = \mathbb{P}\{A^{(k)}_{i} = 1\ |\ B\}
\end{equation}

\section{Methodology}
\label{A5-k-outof-}
In this section I will describe the statistical modeling underlying the approach, and then derive from it the reliability evaluation metrics.  

\subsection{Layer Reliability}
Starting from the assumption that a neural network in operation can be assimilated to the set of the activations of its neurons, then we can state that \textit{the reliability of a layer is a function of the number of activations of its neurons}. More specifically, it is a function of the reliability of a \textit{sub-group} of these in being active when the success occurs.
In reliability theory, this characteristic can be represented by means of a \textit{k-out-of-n} system
\begin{equation}
    \mathcal{R}_{sistema} = \sum_{j=k}^{n}\binom{n}{j} R^{j}(1-R)^{(n-j)} \qquad k=1,\dots, n
\end{equation}
where 
\begin{align}
    &\binom{n}{k} R^{k} \longrightarrow \text{Probability that $k$ components out of $n$ function correctly} \\
    &\binom{n}{k} (1-R)^{(n-k)} \longrightarrow \text{Probability that the remaining $(n-k)$ components out of $n$ do not function}
\end{align}
This formulation uses the \textit{Binomial r.v.} to model the probability of the components. However, for the reliability of a layer of a neural network, it is not possible to use the Binomial since it imposes the \textit{i.i.d.} conditions on the r.v.'s it describes. If, by simplifying the concept of neural network, we could accept the \textit{independence} condition; the \textit{identically distributed} condition cannot be accepted since it would impose on every component of the layer (neuron) to have the \textit{same} reliability distribution.
For this reason, the solution I decided to adopt is to use the generalized version of the Binomial, namely the \textit{Beta-Binomial r.v.}\footnote{In probability theory the Beta-Binomial distribution is a family of discrete probability distributions that can be seen as a generalization of the binomial distribution and of the Beta distribution. It describes the distribution of the number of successes over \textit{n} independent yes/no experiments, but, contrary to the Binomial distribution, the probability of success is not a known parameter, but is an uncertain value distributed as a Beta random variable.}

If $L_{(k)}\sim BetaBin(n_{(k)},\alpha_{(k)},\beta_{(k)})$ is a layer-$k$ defined as a random variable distributed as a Beta-Binomial r.v. with parameters $n_{(k)},\alpha_{(k)},\beta_{(k)}$. Then, with $x \ge 0$ indicating the \textit{number of active components}, we have
\begin{equation}\label{BetaBin}
    \mathbb{P}\{L_{(k)} = x\} = \binom{n_{(k)}}{x}\dfrac{\textbf{B}(\alpha_{(k)}+x\ ,\ \beta_{(k)}+n_{(k)}-x)}{\textbf{B}(\alpha_{(k)}\ ,\ \beta_{(k)})}
\end{equation}
with
\begin{equation}
    \mathbb{E}[X] = n\dfrac{\alpha}{\alpha+\beta} \qquad 
    \mathbb{V}[X] = n\dfrac{\alpha[n(1+\alpha)+\beta]}{(\alpha+\beta)(1+\alpha+\beta)}
\end{equation}
where \textbf{B} is the \textit{Euler Beta} function:
$\textbf{B}(a, b) = \int_{0}^{1} y^{\ a-1}(1-y)^{\ b-1} dy$ \footnote{This distribution governs the probability 
\textit{p} of a Bernoulli process a posteriori of the observation of $(\alpha-1)$ ``successes'' and $(\beta-1)$ ``failures'', when $p$ is a priori uniformly distributed over $[0,1]$}.

The Beta-Binomial handles over-dispersion, that is \textit{high variability}; indeed, while in a classical Binomial the variance is tied to the mean, in a layer it is very easy to have neurons that always activate and others that never do, and this entails a high variance. The Beta-Binomial manages to capture precisely this \textit{heterogeneity} of the neurons within the same layer.

At this point we can define the \textit{layer reliability} as in the \textit{k-out-of-n} model by replacing the Binomial distribution with the Beta-Binomial
\begin{align}\label{Affidabilità Layer}
    \mathcal{R}_{(k)}(n_{(k)},\alpha_{(k)}, \beta_{(k)}) &= \sum_{j=s_{(k)}}^{n_{(k)}} \binom{n_{(k)}}{j}\dfrac{\textbf{B}(\alpha_{(k)}+j\ ,\ \beta_{(k)}+n_{(k)}-j)}{\textbf{B}(\alpha_{(k)}\ ,\ \beta_{(k)})}\\
    &=\sum_{j=s_{(k)}}^{n_{(k)}} \binom{n_{(k)}}{j}\dfrac{\int_{0}^{1} y^{\ (\alpha_{(k)}+j-1)}(1-y)^{\ (\beta_{(k)}+n_{(k)}-j-1)} dy}{\int_{0}^{1} y^{\ \alpha_{(k)}-1}(1-y)^{\ \beta_{(k)}-1} dy}
\end{align}
The reliability of a layer of the network is defined as the \textit{total probability mass} of success, that is, the probability of observing a behavior: \textit{the activation of a sub-group of nodes equal to at least} $s_{(k)}$, which is statistically \textit{significant} when the network returns a correct output.

\subsection{Estimation of the parameters $\alpha$ and $\beta$}
The Beta-Binomial distribution [\ref{BetaBin}] used to describe a layer depends on three parameters: $\alpha_{(k)}, \beta_{(k)}$ and $n_{(k)}$. In order to derive an estimate of the probability distribution of the layer it is therefore necessary to estimate these parameters. More specifically, in our case we need to obtain an estimate of the parameters $\alpha_{(k)}$ and $\beta_{(k)}$, since $n_{(k)}$, being the number of nodes present in the layer, is a known value at the moment the network is built.

To estimate the parameters I used the \textit{Method of Moments} since it is often preferable to maximum likelihood \footnote{Closed-form maximum likelihood estimates are impractical, since the pdf consists of common functions (gamma function and/or beta functions), they can be easily obtained through direct numerical optimization. Maximum likelihood estimates from empirical data can be computed using general methods for fitting multinomial Pólya distributions.} for computational stability and interpretability especially for more complex distributions. The method of moments is based on the equality between the theoretical equation of the \textit{population moments} ($\mu_{r}$) and the \textit{sample} ones ($m_{r}$). In our case, since the parameters to be estimated are two, we will build a system of two equations composed of the equality between the sample mean $m_{1}$ and the population expected value $\mu_{1}$ and the sample variance $m_{2}$ and the population variance $\mu_{2}$.
\begin{equation}
    \begin{cases}
        m_{1} &= \mu_{1}\\
        m_{2} &= \mu_{2}
    \end{cases} \qquad
    \begin{cases}
        \bar{\mathbf{x}}_{n} &= \mathbb{E}[L_{(k)}]\\
        S^{2} &= \mathbb{V}[L_{(k)}]
    \end{cases} \qquad
    \begin{cases}
        \bar{\mathbf{x}}_{n} &= n_{(k)}\dfrac{\alpha_{(k)}}{\alpha_{(k)}+\beta_{(k)}}\\
        S^{2} &= n_{(k)}\dfrac{\alpha_{(k)}[n_{(k)}(1+\alpha_{(k)})+\beta_{(k)}]}{(\alpha_{(k)}+\beta_{(k)})(1+\alpha_{(k)}+\beta_{(k)})}
    \end{cases}
\end{equation}
The first equation is solved for $\alpha_{(k)}$ and we obtain $\frac{\bar{\mathbf{x}}_{n}\beta_{(k)}}{n_{(k)}-\bar{\mathbf{x}}_{n}}$, substituting this value into the second equation and inverting it for $\beta_{(k)}$ we obtain the \textit{point estimates} of the parameters
\begin{equation}\label{stime parametri}
    \begin{cases}
        \hat\alpha_{(k)} &= \dfrac{n_{(k)}m_{1}-m_{2}}{n_{(k)}(\frac{m_{2}}{m_{1}}-{m_{1}} -1)+m_{1}}\\ \\
        \hat\beta_{(k)} &= \dfrac{(n_{(k)}-m_{1})(n_{(k)}-\frac{m_{2}}{m_{1}})}{n_{(k)}(\frac{m_{2}}{m_{1}}-m_{1} -1)+m_{1}}
    \end{cases} \quad \mbox{with } m_{1} = \bar{\mathbf{x}}_{n}\mbox{ and } m_{2} = S^{2}
\end{equation}
Now we can use them to estimate the probability distribution of the Beta-Binomial and use it to study the threshold $s_{(k)}$.\footnote{In the following sections of the work the analysis will be extended by handling the uncertainty of the estimates through the construction of confidence intervals.}

\subsection{Derivation of the threshold $s_{(k)}$}
From the definition of the layer reliability [\ref{Affidabilità Layer}] we said that the parameter $s_{(k)}$ represents the smallest number of neurons that must be active so that the overall activation of the layer is statistically significant with respect to the success behavior of the network. In formal terms we will say that $s_{(k)}$ \textit{is the smallest integer value for which the cumulative probability reaches or exceeds} $q$ (e.g. $0.05$, $0.01$).
\begin{equation}
    s_{(k)} = min_{x}\{ x\in \{0,1, 2, \dots, n_{(k)}\}\ |\ \mathbf{F}_{BetaBin}(x;\ n_{(k)},\alpha_{(k)}, \beta_{(k)})\ge q\}
\end{equation}
Finding this value analytically is very complicated since $\mathbf{F}_{BetaBin}$ is expressed through a \textit{Generalized Hypergeometric} and differentiating it is formally complex. However, we can make an approximation. Indeed, in deep learning the number of neurons per layer is often very large; exploiting this characteristic, we can refer to the asymptotic behavior of the Beta-Binomial distribution. For very large $n_{(k)}$ the Beta-Binomial distribution tends to the Normal, from which we can obtain a \textit{direct} analytical formula for $s_{(k)}$
\begin{equation}\label{soglia sk}
    \lceil s_{(k)}\rceil \approx \hat{\mu}_{(k)} + \Phi^{-1}_q\cdot\hat{\sigma}_{(k)} \qquad \text{with } n_{(k)}\rightarrow \infty
\end{equation}
from which
\begin{align}
    \hat{\mu}_{(k)} &= \mathbb{E}[\hat{L}_{(k)}] = n_{(k)}\dfrac{\hat{\alpha}_{(k)}}{\hat{\alpha}_{(k)}+\hat{\beta}_{(k)}}\\
    \hat{\sigma}_{(k)} &= \sqrt{\mathbb{V}[\hat{L}_{(k)}]} = \sqrt{n_{(k)}\dfrac{\hat{\alpha}_{(k)}[n_{(k)}(1+\hat{\alpha}_{(k)})+\hat{\beta}_{(k)}]}{(\hat{\alpha}_{(k)}+\hat{\beta}_{(k)})(1+\hat{\alpha}_{(k)}+\hat{\beta}_{(k)})}}
\end{align}
$s_{(k)}$ is no longer a theoretical number but rather \textit{the minimum behavior observed under success conditions in layer-k}.

\subsection{Reliability analysis}
Given the definition of reliability [\ref{Affidabilità Layer}], a layer-\textit{k} represented as $L_{(k)}\sim BetaBin(n, \alpha, \beta)$, the parameters $\hat{\alpha}_{(k)}$, $\hat{\beta}_{(k)}$ and the minimum threshold $s_{(k)}$; the \textit{reliability} of the layer is described by the probability:
\begin{equation}
    \mathcal{R}_{(k)}(n_{(k)}, \hat{\alpha}_{(k)}, \hat{\beta}_{(k)}) = \mathbb{P}\{L_{(k)}\ge s_{(k)}\ |\ \hat{\alpha}_{(k)}, \hat{\beta}_{(k)}\} = 1-q
\end{equation}
We are saying that in the analysis of the network we impose \textit{a priori} a reliability value equal to $[(1-q)\cdot 100]\%$ for every layer, each of which will \textit{guarantee such reliability by activating a number of neurons at least equal to} $s_{(k)}$. From here, the reliability analysis focuses \textit{a posteriori} on seeing the number $s_{(k)}$ of neurons that each layer activates in order to guarantee the imposed reliability.

We observe the metric $\rho_{(k)}$ built to represent the behavior of the layer
\begin{equation}
    \rho_{(k)} = \bigg[\dfrac{s_{(k)}}{n_{(k)}}\cdot100\bigg]\%
\end{equation}
This indicates the minimum percentage of active neurons that the layer must have in order to guarantee a reliability equal to $[(1-q)\cdot 100]\%$. For example, if $\rho_{(k)}=10\%$ with $q=0.05$ then we will say that layer-\textit{k}, in order to guarantee a reliability of $95\%$, must activate at least $10\%$ of its neurons.

\textit{How can this measure be interpreted?}\\
A low value of $\rho_{(k)}$ indicates that the layer only needs to activate a small sub-group of neurons in order to be statistically significant at $[(1-q)\cdot100]\%$ when the network achieves a success. It could be interpreted as a signal of robustness of the layer, since it has a lot of information redundancy, on the other hand a layer with a high value of $\rho_{(k)}$ would indicate that the information is dispersed among the neurons. However, tests should be carried out to validate this hypothesis.

\subsection{Confidence Intervals}
The whole approach described so far has used, for the reliability analysis, the point estimates of the parameters $\alpha$ and $\beta$. However, in order to obtain a more robust measure of the estimate of a parameter it is good practice to construct a \textit{confidence interval} \footnote{If $U$ and $V$ are random variables with probability distributions that depend on some parameter $\theta$ and $\mathbb{P}\{U\le \theta \le V\} = 1-\beta$ where $\beta \in [0, 1]$, then the random interval $(U, V)$ computed on the observed sample is a $[(1-\beta)\cdot100]\%$ confidence interval for $\theta$. The extreme values of the confidence interval are called confidence limits.} in which such estimate will move with a certain probability. To construct such an interval one can follow two paths, one analytical defined by a formal procedure called the \textit{Delta Method}. The second one, instead, follows the computational path using the \textit{Bootstrap Method} estimating the intervals from collected samples. In this work, the confidence intervals are derived using the \textit{Delta Method} \footnote{The Delta Method consists of 3 phases: computation of the gradient (or Jacobian matrix) of the function of the parameters, construction of the covariance matrix of the parameters with derivation of the variance of the estimators, and finally construction of the interval around the point estimate of the parameters using the variance.}.\\
This procedure uses the result of the \textit{Central Limit Theorem} on the sample mean $\bar{\mathbf{x}}_{n}$ and the sample variance $S^{2}$, which as $n$ grows tend to the \textit{Multivariate Normal} distribution. The method transfers this normality onto the estimates of the parameters $\hat{\alpha}_{(k)},\hat{\beta}_{(k)}$ using a linearization technique.

\subsubsection{CI construction for $\hat{\alpha}_{(k)}, \hat{\beta}_{(k)}$}
As a first step we construct the \textit{vector function} that we must then use in the procedure.
\begin{equation}
    g(m_{1}, m_{2}) = 
    \begin{bmatrix}
        \hat{\alpha}_{(k)}\\
        \hat{\beta}_{(k)}
    \end{bmatrix}
    =
    \begin{bmatrix}
        \dfrac{n_{(k)}m_{1}-m_{2}}{n_{(k)}(\frac{m_{2}}{m_{1}}-{m_{1}} -1)+m_{1}})\\\\
        \dfrac{(n_{(k)}-m_{1})(n_{(k)}-\frac{m_{2}}{m_{1}})}{n_{(k)}(\frac{m_{2}}{m_{1}}-m_{1} -1)+m_{1}}
    \end{bmatrix} 
    =
    \begin{bmatrix}
        \dfrac{m_{1}(nm_{1}-m_{2})}{m_{1}^{2}-n(m_{1}^{2}+m_{1}-m_{2})})\\\\
        \dfrac{(n-m_{1})(nm_{1}-m_{2})}{(1-n)m_{1}^{2} - nm_{1} + nm_{2}}
    \end{bmatrix}
\end{equation}
Where $m_{1}=\bar{\mathbf{x}}_{n},\ m_{2}=S^{2}$ and the last version is a rearrangement of the terms so as not to have fractions inside the formulas.
Starting from the vector function it is necessary to construct the \textit{Jacobian Matrix} of the partial derivatives.
\begin{equation}
    \mathbf{J}_{2\times2} = \nabla g(m_{1}, m_{2}) = 
    \begin{bmatrix}
        \dfrac{\partial}{\partial m_{1}}\hat{\alpha}_{k} & \dfrac{\partial}{\partial m_{2}}\hat{\alpha}_{k}\\ \\
        \dfrac{\partial}{\partial m_{1}}\hat{\beta}_{k} & \dfrac{\partial}{\partial m_{2}}\hat{\beta}_{k}
    \end{bmatrix}_{2\times2}
\end{equation}
The partial derivatives to be solved for $\hat{\alpha}_{(k)}, \hat{\beta}_{(k)}$ are derivatives of ratios, so the differentiation rule $\Bigl(\frac{\partial}{\partial x}\frac{f(x)}{g(x)} = \frac{f'(x)g(x)-f(x)g'(x)}{(g(x))^{2}} \Bigr)$ is used. Following this procedure and factoring the expressions (in order to have a form as simple as possible), the following Jacobian matrix $\mathbf{J}_{2\times2}$ is obtained. 
\begin{equation}
    \begin{bmatrix}
        \dfrac{m_{1}^{2}(1-n_{(k)})(m_{2}-n_{(k)}^{2}) - n_{(k)}(n_{(k)}m_{1}-b)^{2}}{(m_{1}^{2}(1-n_{(k)}) - n_{(k)}(m_{1}-m_{2}))^{2}} & 
        \dfrac{m_{1}^{2}(n_{(k)}-1)(m_{1}-n_{(k)})}{(m_{1}^{2}(1-n_{(k)}) - n_{(k)}(m_{1}-m_{2}))^{2}} \\ \\
        \dfrac{m_{1}(1-n_{(k)})(m_{1}-n_{(k)})^{2}}{(m_{1}^{2}(1-n_{(k)}) - n_{(k)}(m_{1}-m_{2}))^{2}} & 
        \dfrac{n_{(k)}(n_{(k)}m_{1}-m_{2})^{2} - m_{2}(1-n_{(k)})(m_{1}-n_{(k)})^{2}}{(m_{1}^{2}(1-n_{(k)}) - n_{(k)}(m_{1}-m_{2}))^{2}}
    \end{bmatrix}_{2\times2}
\end{equation}

Once $\mathbf{J}_{2\times2}$ has been obtained it is necessary to construct the \textit{Covariance Matrix} $\large\Sigma_{2\times2}$ of the sample mean and sample variance \footnote{In Beta-Binomial distributions, which are often asymmetric, the term $\hat{\mu}_{3}$ is therefore different from zero. This means that an error in the estimation of the mean propagates causing an error in the estimation of the variance. The matrix captures this relationship, making the final confidence interval much more accurate.}. In the case of the study we are carrying out, the variance-covariance matrix is built on the values observed during the \textit{training} of the network. We construct the vector $\mathbf{x} = \{x_{1}, x_{2}, \dots, x_{N}\}$ of the number of activations \textit{observed} in layer-\textit{k} when the network returns the correct output, with $N$ the total number of correct predictions made during training. The terms of the matrix depend on the \textit{sample central moments} of the distribution of the data $\mathbf{x}$.
\begin{equation}
    \hat{\mu}_{j} = \dfrac{1}{N}\sum_{i=1}^{N}(x_{i}-m_{x})^{j}\qquad\mbox{where } m_{x} = \dfrac{1}{N}\sum_{i=1}^{N}x_{i} \quad\mbox{(sample mean)}
\end{equation}
The matrix $\large\Sigma_{2\times2}$ is populated using precisely the sample central moments normalized by the sample size $N$.
\begin{equation}
    \Sigma_{2\times2} = \dfrac{1}{N} 
    \begin{bmatrix}
        \hat{\mu}_{2} & \hat{\mu}_{3} \\ \\
        \hat{\mu}_{3} & \hat{\mu}_{4} - (\hat{\mu}_{2})^{2}
    \end{bmatrix}_{2\times2}
\end{equation}
In order to construct the confidence interval around the point estimates $\hat{\alpha}_{k}, \hat{\beta}_{k}$ it is necessary to obtain the \textit{variance of the estimators}. These values are obtained by carrying out the \textit{sandwich} matrix multiplication between the matrices $\mathbf{J}_{2\times2}$ and $\large\Sigma_{2\times2}$ and taking the values on the main diagonal of the resulting matrix.
\begin{equation}\label{matriceSandwich}
     \mathbf{J}_{2\times2}\ \Sigma_{2\times2}\ (\mathbf{J}_{2\times2})^{T} =
     \begin{bmatrix}
         \mathbb{V}(\hat{\alpha}_{(k)}) & \text{Cov}(\hat{\alpha}_{(k)}, \hat{\beta}_{(k)}) \\ \\
         \text{Cov}(\hat{\alpha}_{(k)}, \hat{\beta}_{(k)}) & \mathbb{V}(\hat{\beta}_{(k)})
     \end{bmatrix}_{2\times2}
\end{equation}
the matrix obtained is nothing other than the \textit{variance-covariance} matrix of the estimators $\hat{\alpha}_{k}, \hat{\beta}_{k}$. 

Now we can construct the asymptotic confidence intervals for the point estimates by letting them vary in a way that depends on their estimated variance.
\begin{equation}
    CI_{q}[\alpha_{(k)}]
    = \hat{\alpha}_{(k)}\pm \Phi^{-1}_q\sqrt{\mathbb{V}(\hat{\alpha}_{(k)})} \qquad \quad
    CI_{q}[\beta_{(k)}]
    = \hat{\beta}_{(k)}\pm \Phi^{-1}_q\sqrt{\mathbb{V}(\hat{\beta}_{(k)})}
\end{equation}
where $\mathbb{V}(\hat{\alpha}_{(k)}) = \Bigl[\mathbf{J}_{2\times2}\ \Sigma_{2\times2}\ (\mathbf{J}_{2\times2})^{T}\Bigr]_{1,1}$ and $\mathbb{V}(\hat{\beta}_{(k)}) = \Bigl[\mathbf{J}_{2\times2}\ \Sigma_{2\times2}\ (\mathbf{J}_{2\times2})^{T}\Bigr]_{2,2}$.

\subsubsection{CI construction for $R_{(k)}$}
As seen in the analysis [\ref{soglia sk}], we studied the reliability $\mathcal{R}_{(k)}$ of a layer by observing the estimate of the minimum number of neurons that must be active $s_{(k)}$ in order to guarantee an a priori imposed reliability. However, we know that this also depends on the point estimates $\hat{\alpha}_{(k)}, \hat{\beta}_{(k)}$ [\ref{Affidabilità Layer}], so it is necessary to study in which interval $\mathcal{R}_{(k)}$ moves as $\hat{\alpha}_{(k)}, \hat{\beta}_{(k)}$ vary.

As a first step we obtain the \textit{Sensitivity} of $\mathcal{R}_{(k)}$ that is $\nabla \mathcal{R}_{(k)}(n_{(k)}, \alpha_{(k)}, \beta_{(k)})$ which tells us how $\mathcal{R}_{(k)}$ moves as the parameters $\alpha_{(k)}, \beta_{(k)}$ vary.
\begin{equation} \label{vettoreSensibilità}
    \mathbf{R}_{1\times2} = \nabla \mathcal{R}_{(k)}(n_{(k)}, \alpha_{(k)}, \beta_{(k)}) = 
    \begin{bmatrix}
        \dfrac{\partial}{\partial \alpha} \mathcal{R}_{(k)}(n_{(k)}, \alpha_{(k)}, \beta_{(k)}) & \dfrac{\partial}{\partial \beta} \mathcal{R}_{(k)}(n_{(k)}, \alpha_{(k)}, \beta_{(k)})
    \end{bmatrix}
\end{equation}
We consider the reliability formula $\mathcal{R}_{(k)}(n_{(k)}, \alpha_{(k)}, \beta_{(k)})$ [\ref{Affidabilità Layer}], by developing the integrals using the rule $\int_{0}^{1}x^{S}dx = \frac{1}{S+1}$, carrying out the computations and rearranging the terms we obtain a new form of $\mathcal{R}_{(k)}$ that is easier to handle for constructing $\mathbf{R}_{1\times2}$.
\begin{equation}
    \mathcal{R}_{(k)}(n_{(k)}, \alpha_{(k)}, \beta_{(k)}) = \sum_{j = s_{(k)}}^{n_{(k)}} \binom{n_{(k)}}{j} 
    \dfrac{\frac{1}{A(\alpha_{(k)})} - \frac{1}{B(\alpha_{(k)}, \beta_{(k)})}}{\frac{1}{C(\alpha_{(k)})} - \frac{1}{D(\alpha_{(k)}, \beta_{(k)})}}
\end{equation}
where $A(\alpha_{(k)}) = \alpha_{(k)} + j,
    B(\alpha_{(k)}, \beta_{(k)}) = \alpha_{(k)} + \beta_{(k)} + n - 1,
    C(\alpha_{(k)}) = \alpha_{(k)} \text{ and }
    D(\alpha_{(k)}, \beta_{(k)}) = \alpha_{(k)} + \beta_{(k)} - 1$.
Using this formulation, although it turns out to be less complex than the original form, in order to obtain the simplest possible partial derivatives with respect to $\alpha_{(k)}, \beta_{(k)}$, it is necessary to rearrange the expression in the form $Kf(x)$ so as to have a constant term $K$ that does not depend on the parameter with respect to which we are differentiating. In this way, both the computations and the final results are in the simplest possible form.

We start from the derivative of $\mathcal{R}_{(k)}(n_{(k)}, \alpha_{(k)}, \beta_{(k)})$ with respect to $\alpha_{(k)}$.\\
The reliability expressed in $Kf(x)$ form is the following:
\begin{equation}
    \mathcal{R}_{(k)}(n_{(k)}, \alpha_{(k)}, \beta_{(k)}) = \sum_{j = s_{(k)}}^{n_{(k)}} \binom{n_{(k)}}{j} K_{1}f_{1}(\alpha_{(k)})
\end{equation}
with
\begin{equation}
    K_{1} = \dfrac{\beta_{(k)} + n_{(k)} - j - 1}{\beta_{(k)} - 1} \qquad
    f_{1}(\alpha_{(k)}) = \dfrac{\alpha_{(k)}(\alpha_{(k)}+\beta_{(k)} -1)}{(\alpha_{(k)} + j)(\alpha_{(k)}+\beta_{(k)}+n-1)}
\end{equation}
We compute the partial derivative with respect to $\alpha_{(k)}$. Instead of differentiating using the classical quotient rule (it would turn out very complex), I use the logarithm rule $\frac{\partial}{\partial x}f(x) = f(x)\frac{\partial}{\partial x}\ln(f(x))$, which turns out to be simpler to differentiate.
\begin{equation}
    \dfrac{\partial}{\partial \alpha_{(k)}}\mathcal{R}_{(k)}= \sum_{j = s_{(k)}}^{n_{(k)}} \binom{n_{(k)}}{j} K_{1}f_{1}(\alpha_{(k)})
    \Bigl[ \frac{1}{\alpha_{(k)}} + \frac{1}{\alpha_{(k)} + \beta_{(k)} -1} -\frac{1}{\alpha_{(k)}+j} - \frac{1}{\alpha_{(k)} + \beta_{(k)} + n -1}\Bigr]
\end{equation}

We move on to the derivative of $\mathcal{R}_{(k)}(n_{(k)}, \alpha_{(k)}, \beta_{(k)})$ with respect to $\beta_{(k)}$.\\
The reliability expressed in $Kf(x)$ form is the following:
\begin{equation}
    \mathcal{R}_{(k)}(n_{(k)}, \alpha_{(k)}, \beta_{(k)}) = \sum_{j = s_{(k)}}^{n_{(k)}} \binom{n_{(k)}}{j} K_{2}f_{2}(\beta_{(k)})
\end{equation}
With
\begin{equation}
    K_{2} = \dfrac{\alpha_{(k)}}{\alpha_{(k)} + j} \qquad
    f_{2}(\beta_{(k)}) = \dfrac{(\alpha_{(k)} + \beta_{(k)} - 1)(\beta_{(k)} + n - j - 1)}{(\beta_{(k)}-1)(\alpha_{(k)}+ \beta_{(k)} + n -1)}
\end{equation}
We compute the partial derivative with respect to $\beta_{(k)}$, here too I use the logarithm rule to differentiate the function.
\begin{equation}
    \dfrac{\partial}{\partial \beta_{(k)}}\mathcal{R}_{(k)}= \sum_{j = s_{(k)}}^{n_{(k)}} \binom{n_{(k)}}{j} K_{2}f_{2}(\beta_{(k)})
    \Bigl[ \frac{1}{\alpha_{(k)} + \beta_{(k)} - 1} + \frac{1}{\beta_{(k)} + n - j - 1} -\frac{1}{\alpha_{(k)}+\beta_{(k)} + n -1} - \frac{1}{\beta_{(k)} - 1}\Bigr]
\end{equation}
We have obtained the \textit{sensitivity} vector $\mathbf{R}_{1\times2}$ [\ref{vettoreSensibilità}] of $\mathcal{R}_{(k)}$ with respect to $\alpha_{(k)}, \beta_{(k)}$.

Following the Delta method, one must carry out the \textit{sandwich} matrix multiplication between $\mathbf{R}_{1\times2}$ and the covariance matrix of the estimates of the parameters $\alpha_{(k)}, \beta_{(k)}$. But we have already obtained this matrix during the construction of the confidence interval of the parameters when we carried out the matrix multiplication $\mathbf{J}_{2\times2}\ \Sigma_{2\times2}\ (\mathbf{J}_{2\times2})^{T}$ [\ref{matriceSandwich}]. Denoting this matrix with $\hat{\Sigma}_{2\times2}$, we can obtain the \textit{variance} of $\mathcal{R}_{(k)}(n_{(k)}, \alpha_{(k)}, \beta_{(k)})$.
\begin{equation}
    \mathbb{V}[\mathcal{R}_{(k)}(n_{(k)}, \alpha_{(k)}, \beta_{(k)})] \approx (\mathbf{R}_{1\times2})^{T}\ \hat{\Sigma}_{2\times2} \ (\mathbf{R}_{1\times2})  
\end{equation}
The result is a scalar that represents the variance of the reliability $\mathcal{R}_{(k)}(n_{(k)}, \alpha_{(k)}, \beta_{(k)})$ of layer-\textit{k}. Once this value has been obtained it is possible to construct the \textit{confidence interval} of the reliability around its estimate at the points $\hat{\alpha}_{(k)},\hat{\beta}_{(k)}$.
\begin{equation}
    CI_{q}[\mathcal{R}_{(k)}(n_{(k)}, \alpha_{(k)}, \beta_{(k)})] = \mathcal{R}_{(k)}(n_{(k)}, \hat{\alpha}_{(k)}, \hat{\beta}_{(k)}) \pm \Phi^{-1}_{q}\sqrt{\mathbb{V}[\mathcal{R}_{(k)}(n_{(k)}, \alpha_{(k)}, \beta_{(k)})]}
\end{equation}
This confidence interval allows us to quantify the \textit{uncertainty} associated with the point estimate $\mathcal{R}_{(k)}$, so as to be able to carry out a critical evaluation of neural networks as is done in \textit{safety-critical} contexts for classical systems.
\begin{enumerate}
    \item Identification of the \textit{Lower Confidence Bound}. \\
    In the safety-critical analysis of systems one is interested in the \textit{worst case}. In this case, the lower endpoint of the interval represents the \textit{minimum guaranteed reliability} with a certain specified confidence level (e.g. $95\%$). Therefore, once the interval has been obtained, the lower limit allows us to state that: \textit{Layer-k, with a minimum activation threshold of neurons equal to $s_{(k)}$, guarantees a reliability at least equal to the lower endpoint of the interval with a confidence of $[(1-q)\cdot100]\%$}\footnote{Suppose that after training, following the procedure described so far, $\hat{\alpha}_{(k)}, \hat{\beta}_{(k)}$ are estimated and the metric $\rho_{(k)} = [s_{(k)}/{n_{(k)}}\cdot100]\% = 40\%$ is obtained for layer-\textit{k}. By constructing the confidence interval around the reliability value $R_{(k)}$ at $\hat{\alpha}_{(k)}, \hat{\beta}_{(k)}$ with a confidence of $95\%$, a lower limit of approximately $75\%$ is observed. Then we will say that layer-\textit{k} will have a reliability of at least $75\%$, $95\%$ of the time, by activating about $40\%$ of its neurons.}.
    
    \item \textit{Training robustness} metric.\footnote{The study of how the interval behaves asymptotically and as a function of which values will be addressed in future works.}\\
    The width of the interval could also be seen as a measure related to the quality of the training for a specific layer.
    \begin{itemize}
        \item A narrow interval implies \textit{less} variability and therefore a \textit{stable} convergence;
        \item A wide interval indicates that, regardless of accuracy, the behavior inside the layer is \textit{volatile}.
    \end{itemize}
\end{enumerate}

\section{Conclusions}
The study is an attempt to introduce a method for analyzing neural networks, proposing a shift from the simple performance-based evaluation to a \textit{reliability-based} evaluation. By modeling the layers as \textit{k}-out-of-\textit{n} stochastic systems represented by Beta-Binomial distributions, it is possible to quantify the \textit{participation of the neurons} and to define data-driven activation thresholds such as $s_{(k)}$ and operational metrics such as $\rho_{(k)}$.

% However, the most significant contribution of this approach, in my opinion, lies in its use for Explainable AI (XAI). Since traditional XAI techniques focus on analyzing \textit{why} a network made a certain decision, complementing them with this approach means also answering the question: \textit{With how much internal reliability is the network operating?}
 
% Since this work focused on the definition and the theoretical demonstration of the methodology, the next steps I set for myself are
% \begin{itemize}
%     \item Broaden the theoretical study, by studying the functions that have been constructed for the reliability and the confidence intervals, in order to deduce other structural characteristics from them;
%     \item Development of an application framework capable of \textit{automating the statistical inference}, \textit{runtime monitoring} of the reliability of the layers and their \textit{validation}.
% \end{itemize}

\end{document}
